\section{Análisis transitorio}

La primera simulación corresponde al circuito RLC de la consigna. La fuente conmuta entre \SI{0}{V} y \SI{10}{V}, y se observa la tensión $v_o(t)$ sobre la rama del inductor.

\subsection{Circuito analizado}

\begin{figure}[H]
    \centering
    \begin{tikzpicture}[transform shape]
        % Fuente de entrada con signos + y - (v={$V_i$})
        \draw (5, 1.5) to[square voltage source, v_={$V_i$}, american voltages] (5, -1.5);
        
        % Resistencias superiores
        \draw (6.25, 1.5) to[american resistor, l={$R_i$}] (7.25, 1.5);
        \draw (8.25, 1.5) to[american resistor, l={$R_1$}] (9.5, 1.5);
        
        % Rama del capacitor y R2
        \draw (11.5, 1.25) to[capacitor, l={$C$}] (11.5, 0.25);
        \draw (11.5, 0.25) to[american resistor, l={$R_2$}] (11.5, -1.5);
        
        % Rama del inductor y R3
        \draw (13, 1.5) to[american inductor, l={$L$}] (13, -0);
        \draw (13, -0) to[american resistor, l={$R_3$}] (13, -1.25);
        
        % Conexiones de cables (líneas horizontales y verticales)
        \draw (5, 1.5) -- (6.25, 1.5);
        \draw (7.25, 1.5) -- (8.25, 1.5);
        \draw (9.5, 1.5) -- (11.5, 1.5);
        \draw (11.5, 1.5) -| (11.5, 1.25);
        \draw (11.5, 1.5) -- (13, 1.5);
        \draw (5, -1.5) -- (11.5, -1.5);
        \draw (11.5, -1.5) -- (13, -1.5);
        \draw (13, -1.25) -| (13, -1.5);
    
        % Extensión a la derecha para indicar Vo con sus signos + y - en paralelo
        \draw (13, 1.5) -- (14.5, 1.5) 
              to[open, v={$V_o$}, american voltages] (14.5, -1.5) 
              -- (13, -1.5);
    \end{tikzpicture}
    \caption{Circuito utilizado para el análisis transitorio.}
    \label{fig:circuito_transitorio}
\end{figure}

Los valores utilizados fueron:

\begin{equation}
    V_i=\SI{10}{V}, \qquad L=\SI{1}{mH}, \qquad R_1=\SI{50}{\ohm},
\end{equation}

\begin{equation}
    R_2=\SI{8.2}{\ohm}, \qquad R_3=\SI{6.8}{\ohm}, \qquad C=\SI{47}{nF}.
\end{equation}

En la simulación se agregó además la resistencia interna del generador y resistencias parásitas de los componentes cuando fue necesario, ya que influyen directamente en el amortiguamiento.

\subsection{Modelo teórico}

Durante cada intervalo en el que la llave permanece fija, el circuito se puede describir como un sistema lineal de segundo orden. La respuesta natural de una variable del circuito tiene la forma:

\begin{equation}
    x(t)=x_{\infty}+e^{-\alpha t}\left(A\cos(\omega_d t)+B\sin(\omega_d t)\right),
\end{equation}

si el caso es subamortiguado. La frecuencia amortiguada es:

\begin{equation}
    \omega_d=\sqrt{\omega_0^2-\alpha^2}.
\end{equation}

La pseudofrecuencia observada en el osciloscopio se obtiene a partir del período entre máximos sucesivos:

\begin{equation}
    f_d=\frac{1}{T_d}.
\end{equation}

El sobrepico se calcula comparando el máximo de la señal con el valor final:

\begin{equation}
    M_p=\frac{V_{\max}-V_f}{|V_f|}\cdot 100\%.
\end{equation}

\subsection{Comparación entre simulación y medición}

Los archivos medidos y simulados no poseen el mismo origen temporal. Por lo tanto, antes de compararlos se alineó cada señal con:

\begin{equation}
    t_{\mathrm{alineado}}=t-t_0.
\end{equation}

Luego se aplicó un desplazamiento horizontal adicional para que los primeros máximos positivos coincidieran. Este procedimiento no modifica la forma ni la amplitud de las señales, solo permite comparar ambas respuestas en fase.

\begin{figure}[H]
    \centering
    \includegraphics[width=0.92\textwidth]{\figpath transitorio_2_1.pdf}
    \caption{Comparación entre transitorio medido y transitorio simulado, alineados temporalmente.}
    \label{fig:transitorio_21}
\end{figure}

Se observa que la simulación reproduce correctamente la forma general del transitorio. Las diferencias se atribuyen a tolerancias de los componentes, resistencia interna del generador, resistencia serie de la bobina, contactos de protoboard y ruido de medición. La medición presenta ruido superpuesto, mientras que la simulación ideal es suave.

\subsection{Caso ideal con resistencias nulas}

Si las resistencias fueran nulas, el circuito no tendría disipación de energía. En un modelo ideal, la energía intercambiaría indefinidamente entre el capacitor y el inductor, obteniéndose una oscilación sin amortiguamiento. En la realidad esto no ocurre porque siempre existen resistencias parásitas: resistencia de cables, resistencia serie de inductores, resistencia interna de la fuente, pérdidas dieléctricas y pérdidas por contacto. Por lo tanto, aunque no se agreguen resistencias externas, el circuito real siempre presenta amortiguamiento.
